\subsection{Witness Cost for Class Identity}

Recall from Section 2 that the witness cost $W(P)$ is the minimum number of primitive queries required to compute property $P$. For class identity, we ask: what is the minimum number of queries to determine if two values have the same class?

\begin{theorem}[Nominal-Tag Observers Achieve Minimum Witness Cost]\label{thm:nominal-minimum-witness-cost}
Nominal-tag observers achieve the minimum witness cost for class identity:
$$W_{\text{eq}} = O(1)$$

Specifically, the witness is a single tag read: compare $\text{tag}(v_1) = \text{tag}(v_2)$.

Attribute-only observers require $W_{\text{eq}} = \Omega(d)$ where $d$ is the distinguishing dimension (and $d \le n$, with worst-case $d = n$).
\end{theorem}
\leanmeta{ \LH{ACS7}, \LH{ACS1}}

\begin{proof}
\begin{enumerate}
\item Nominal-tag access is a single primitive query.
\item Attribute-only observers must query at least $d$ attributes in the worst case (a generic strategy queries all $n$).
\item No shorter witness exists for attribute-only observers by the information barrier.
\end{enumerate}
\end{proof}

\subsection{Witness Cost Comparison}

\begin{table}[h]
\centering
\begin{tabular}{|l|c|c|}
\hline
\textbf{Observer Class} & \textbf{Witness Procedure} & \textbf{Witness Cost $W$} \\
\hline
Nominal-tag & Single tag read & $O(1)$ \\
Attribute-only & Query a distinguishing set & $\Omega(d)$ \\
\hline
\end{tabular}
\vspace{3pt}
\caption{Witness cost for class identity by observer class.}
\end{table}

The supplementary artifact provides a machine-checked proof that nominal-tag access minimizes witness cost for class identity.
